\documentclass[11pt]{article}

\usepackage[margin=1in]{geometry}
\usepackage{amsmath,amssymb}
\usepackage{array}
\usepackage{enumitem}

\title{Lecture 1 Notes: \\ What Is So Interesting About Mathematics for a Philosopher?}
\author{Philosophy of Mathematics}
\date{}

\begin{document}

\maketitle

\section*{Central Theme}

Mathematics is philosophically interesting because it appears to be abstract, necessary,
and knowable by reason alone, while also being indispensable to empirical science and
ordinary inquiry.

\[
\text{Mathematics} = \text{abstract reason} + \text{scientific applicability}
\]

The guiding question for the lecture:

\[
\boxed{\text{Who has authority over mathematics?}}
\]

Philosophers? Mathematicians? Science? Mathematical practice itself?

\section*{1. Mathematics as a Philosophical Test Case}

Mathematics has attracted philosophers throughout history. Plato treated mathematics as
training for grasping reality as it is, rather than merely as it appears. Many major figures
in philosophy were also deeply engaged with mathematics or logic: Descartes, Leibniz,
Frege, Russell, Hilbert, G\"odel, Tarski, and others.

Mathematics is a test case for several major philosophical problems:

\[
\begin{array}{rcl}
\text{Knowledge} &:& \text{How do we know mathematical truths?} \\
\text{Reality} &:& \text{What are numbers, sets, functions, spaces?} \\
\text{Language} &:& \text{What do mathematical terms refer to?} \\
\text{Logic} &:& \text{What counts as valid reasoning?} \\
\text{Science} &:& \text{Why is mathematics so useful in empirical inquiry?}
\end{array}
\]

\subsection*{Teaching Point}

Philosophy of mathematics is not just ``philosophers talking about numbers.'' It is a way
into central questions about knowledge, reality, language, logic, and science.

\subsection*{Whiteboard}

\[
\text{Mathematics}
\quad \longrightarrow \quad
\begin{cases}
\text{Knowledge} \\
\text{Reality} \\
\text{Language} \\
\text{Logic} \\
\text{Science}
\end{cases}
\]

\section*{2. Rationalism and Empiricism}

Mathematics seems to provide knowledge that is certain, necessary, and independent of
sense experience. This makes it central to the dispute between rationalism and empiricism.

\[
\begin{array}{c|c}
\textbf{Rationalism} & \textbf{Empiricism} \\
\hline
\text{Reason is a source of knowledge} & \text{Experience is the source of knowledge} \\
\text{Mathematics is a model} & \text{Mathematics is a problem} \\
\text{Plato, Descartes, Leibniz} & \text{Aristotle, Locke, Hume, Mill}
\end{array}
\]

For the rationalist, mathematics shows what genuine knowledge looks like: proof,
necessity, certainty. For the empiricist, mathematics is a challenge: if knowledge comes
from experience, how do we explain mathematical knowledge?

\subsection*{Teaching Point}

Mathematics appears to be a paradigm case of \emph{a priori} knowledge: knowledge prior
to, or independent of, experience.

\subsection*{Whiteboard}

\[
\text{Mathematics} = \text{paradigm of } a priori \text{ knowledge?}
\]

\[
2+3=5
\qquad
\text{known by proof/reason, not experiment?}
\]

\section*{3. Mathematics and the Sciences}

Mathematics is not merely abstract. It is also indispensable to the natural and social
sciences. Physics, engineering, economics, computer science, statistics, and many parts of
biology depend on mathematical representation.

This creates a central puzzle:

\[
\boxed{\text{Why does abstract mathematics apply to the physical world?}}
\]

Galileo's metaphor is useful here: the book of nature is written in the language of
mathematics.

\subsection*{Teaching Point}

Mathematics seems independent of experience, yet it is one of our most powerful tools for
understanding experience. This tension is one of the deepest motivations for philosophy of
mathematics.

\subsection*{Whiteboard}

\[
\text{Abstract mathematics}
\quad \longrightarrow \quad
\text{empirical science}
\]

\[
\text{Puzzle: Why does mathematics apply?}
\]

\section*{4. Logic as a Bridge Between Mathematics and Philosophy}

Modern formal logic developed largely through reflection on mathematical reasoning.
Figures such as Boole, Frege, Russell, Hilbert, G\"odel, Tarski, and Church were interested
in the foundations of mathematics. But the tools they developed became central to
analytic philosophy more generally.

Logic connects mathematics to questions about proof, truth, consequence, validity,
meaning, and reference.

\subsection*{Teaching Point}

Mathematics did not merely give philosophers something to study. It helped shape the
formal tools philosophers use to study many other things.

\subsection*{Whiteboard}

\[
\text{Mathematics}
\quad \longrightarrow \quad
\text{Logic}
\quad \longrightarrow \quad
\text{Analytic Philosophy}
\]

\[
\text{proof} \qquad
\text{truth} \qquad
\text{validity} \qquad
\text{model} \qquad
\text{meaning}
\]

\section*{5. Philosophy and Mathematics: Chicken or Egg?}

The central methodological question of the chapter is whether philosophy should govern
mathematics, or whether mathematics should govern philosophy.

\[
\boxed{\text{Does philosophy determine mathematics, or does mathematics determine philosophy?}}
\]

The chapter presents three broad positions:

\[
\begin{array}{c|c|c}
\textbf{Philosophy-first} & \textbf{Middle view} & \textbf{Philosophy-last} \\
\hline
\text{Philosophy governs} & \text{Philosophy interprets} & \text{Mathematics governs} \\
\text{First principles} & \text{Mutual illumination} & \text{Practice first} \\
\text{Revisionary} & \text{No domination} & \text{Anti-revisionist}
\end{array}
\]

\subsection*{Teaching Point}

The chapter is not only about mathematics. It is also about disciplinary authority. Who gets
to say what counts as legitimate mathematics?

\subsection*{Whiteboard}

\[
\text{Philosophy-first}
\quad \longleftrightarrow \quad
\text{Philosophy-in-between}
\quad \longleftrightarrow \quad
\text{Philosophy-last-if-at-all}
\]

\section*{6. Plato: Philosophy Pressuring Geometry}

Plato held that the true objects of mathematics are eternal, unchanging, and ideal.
Geometry is not really about the imperfect figures we draw. It is about ideal objects
grasped by reason.

But geometers often use constructive language:

\[
\text{draw a line}
\qquad
\text{move a figure}
\qquad
\text{construct a triangle}
\]

Plato objects that this language is misleading. If mathematical objects are eternal and
unchanging, then they are not literally drawn, moved, or constructed.

\subsection*{Teaching Point}

Plato is an example of philosophy-first. His metaphysics of mathematical objects leads him
to criticize ordinary mathematical language.

\subsection*{Whiteboard}

\[
\textbf{Plato}
\]

\[
\text{Mathematical objects are eternal and unchanging.}
\]

\[
\therefore \quad \text{``draw a line'' is philosophically misleading.}
\]

\section*{7. Intuitionism: Philosophy Pressuring Logic}

Intuitionists such as Brouwer and Heyting hold that mathematics concerns mental
construction. To say that a mathematical object exists is to say, roughly, that it can be
constructed.

This puts pressure on classical logic, especially the law of excluded middle:

\[
P \lor \neg P
\]

Classically, the following inference is valid:

\[
\neg \forall x P(x)
\quad \Rightarrow \quad
\exists x \neg P(x)
\]

But the intuitionist rejects this as a general principle. To assert an existential claim, one
must be able to construct an example.

\subsection*{Teaching Point}

For intuitionists, logic depends on the meaning of mathematical existence. If existence
means constructibility, then some classical principles become suspect.

\subsection*{Whiteboard}

\[
\textbf{Intuitionism}
\]

\[
\text{Mathematics} = \text{mental construction}
\]

\[
\text{To exist} = \text{to be constructible}
\]

\[
\neg \forall x P(x)
\quad \not\Rightarrow_{\text{intuitionist}} \quad
\exists x \neg P(x)
\]

\section*{8. Impredicative Definitions: Poincar\'e and G\"odel}

A definition is impredicative when it defines an object by reference to a totality that already
contains that object.

\[
\text{Impredicative definition}
=
\text{definition using a totality containing the thing defined}
\]

Poincar\'e objected that such definitions are circular. If mathematical objects are
constructed, then one cannot construct an object by appealing to a totality that already
contains it.

G\"odel defended impredicative definitions from a realist perspective. If mathematical
objects exist independently of us, then definitions do not create them. They describe or
identify objects that already exist.

\subsection*{Teaching Point}

Whether an impredicative definition is legitimate depends partly on whether definitions are
understood as recipes for constructing objects or descriptions of already existing objects.

\subsection*{Whiteboard}

\[
\begin{array}{c|c}
\textbf{Poincar\'e} & \textbf{G\"odel} \\
\hline
\text{Objects constructed} & \text{Objects discovered} \\
\text{Definitions as recipes} & \text{Definitions as descriptions} \\
\text{Impredicativity circular} & \text{Impredicativity legitimate}
\end{array}
\]

\section*{9. Anti-Revisionism: Mathematical Practice Pushes Back}

The chapter argues that philosophy-first does not fit the history of mathematics very well.

Classical logic and impredicative definitions are deeply entrenched in modern mathematics.
But mathematicians did not generally accept them because they first adopted a realist
metaphysics. Rather, these methods became accepted because they were mathematically
fruitful.

\[
\text{Classical logic}
+
\text{impredicative definitions}
=
\text{central to modern mathematics}
\]

\subsection*{Teaching Point}

Anti-revisionism says that successful mathematical practice has default authority. A
philosophical argument must be very strong before it can overturn established mathematics.

\subsection*{Whiteboard}

\[
\textbf{Anti-revisionism}
\]

\[
\text{Successful mathematical practice has default authority.}
\]

\[
\text{If philosophy conflicts with successful mathematics,}
\]

\[
\text{maybe philosophy is the problem.}
\]

\section*{10. Dummett: Language and Mathematical Revision}

Michael Dummett gives another example of a philosophical argument with revisionary
consequences. He argues from considerations about language, meaning, learnability, and
communication.

If meaning must be connected to use and verification, then classical mathematics may be
problematic when it appeals to truths that outrun our capacity to recognize them.

This pushes Dummett toward intuitionistic logic.

\subsection*{Teaching Point}

Dummett raises the question of whether a theory of meaning can constrain mathematical
logic.

\subsection*{Whiteboard}

\[
\textbf{Dummett}
\]

\[
\text{Meaning}
\quad \longrightarrow \quad
\text{use, learnability, communication}
\]

\[
\therefore \quad
\text{classical logic is not generally valid?}
\]

\[
\text{Revise mathematics}
\quad \text{or}
\quad
\text{reject the philosophy of language?}
\]

\section*{11. Naturalism: No First Philosophy}

The final section introduces naturalism. Naturalism rejects the idea that philosophy stands
outside science and mathematics as a superior tribunal.

For Quine, philosophy begins within our inherited scientific worldview. We do not justify
science or mathematics from an external standpoint. We revise our beliefs from within our
overall web of belief.

The image is Neurath's boat: we rebuild the ship while already at sea.

\subsection*{Teaching Point}

Naturalism rejects philosophy-first. Philosophy is not prior to inquiry. It is continuous with
inquiry.

\subsection*{Whiteboard}

\[
\textbf{Naturalism}
\]

\[
\text{No first philosophy}
\]

\[
\text{Philosophy is continuous with inquiry}
\]

\[
\text{Neurath's boat: repair the ship while at sea}
\]

\section*{12. Quine and Maddy}

Quine and Maddy both reject philosophy-first, but they differ over the status of
mathematics.

For Quine, mathematics is accepted because of its role in empirical science. Mathematics is
part of the scientific web of belief.

For Maddy, mathematics has its own methods, standards, and internal goals.
Mathematics should not need external justification from either philosophy or immediate
scientific application.

\[
\begin{array}{c|c}
\textbf{Quine} & \textbf{Maddy} \\
\hline
\text{Science-centered naturalism} & \text{Mathematics-centered naturalism} \\
\text{Mathematics justified by science} & \text{Mathematics has its own standards} \\
\text{Web of belief} & \text{Seams in the web} \\
\text{Empirical role matters} & \text{Mathematical practice matters}
\end{array}
\]

\subsection*{Teaching Point}

Quine rejects philosophy-first, but may still pressure mathematics from the standpoint of
science. Maddy gives more autonomy to mathematics itself.

\subsection*{Whiteboard}

\[
\text{Quine: mathematics earns its place through science.}
\]

\[
\text{Maddy: mathematics has autonomous standards.}
\]

\section*{Conclusion: Why Mathematics Matters Philosophically}

Mathematics is philosophically important because it combines features that are difficult to
hold together.

\[
\begin{array}{c|c}
\textbf{Feature of Mathematics} & \textbf{Philosophical Problem} \\
\hline
\text{Abstract} & \text{What is mathematics about?} \\
\text{Necessary} & \text{How can we know necessary truths?} \\
\text{A priori} & \text{How does this fit with empiricism?} \\
\text{Logical} & \text{What principles of reasoning are legitimate?} \\
\text{Applicable} & \text{Why does mathematics describe the world?} \\
\text{Autonomous} & \text{Can philosophy criticize mathematical practice?}
\end{array}
\]

\section*{Final Framing}

To teach philosophy of mathematics well, we should present mathematics not merely as a
body of theorems, but as a philosophical pressure point.

Mathematics appears to be the purest exercise of reason, yet it is indispensable to our
empirical knowledge of the world. It seems autonomous, yet it raises questions that reach
into epistemology, metaphysics, logic, language, and philosophy of science.

\[
\boxed{
\text{Philosophy of mathematics interprets mathematics and its place in inquiry.}
}
\]

\section*{Discussion Questions}

\begin{enumerate}[leftmargin=2em]
    \item Why has mathematics been so attractive to rationalist philosophers?
    \item Why is mathematics difficult for empiricism?
    \item Should philosophers ever tell mathematicians to change their methods?
    \item Does mathematical existence require construction?
    \item Are mathematical definitions recipes or descriptions?
    \item Is mathematics justified by its role in science?
    \item Does pure mathematics need external justification?
    \item Is philosophy of mathematics supposed to ground mathematics, criticize it, or interpret it?
\end{enumerate}

\section*{Key Terms}

\begin{description}[leftmargin=2em]
    \item[Rationalism] The view that reason is a fundamental source of knowledge.

    \item[Empiricism] The view that sense experience is the primary source of knowledge.

    \item[A priori knowledge] Knowledge independent of sense experience.

    \item[Philosophy-first] The view that philosophy supplies first principles that govern mathematical practice.

    \item[Philosophy-last-if-at-all] The view that mathematics is autonomous and philosophy merely describes it afterward, if it has any role at all.

    \item[Anti-revisionism] The view that successful mathematical practice should not be revised merely because of external philosophical arguments.

    \item[Intuitionism] The view that mathematics concerns mental construction and that existence requires constructibility.

    \item[Law of excluded middle] The classical principle that for any proposition $P$, either $P$ or $\neg P$ holds.

    \item[Impredicative definition] A definition that characterizes an object by reference to a totality containing that object.

    \item[Realism] The view that mathematical objects exist independently of minds, language, or constructions.

    \item[Naturalism] The view that philosophy is continuous with successful inquiry rather than an external foundation for it.
\end{description}

\end{document}